  % ── Ⅰ. 서 론 ────────────────────────────────────────────
  \chapter{서\quad 론}
  \section{서론의 내용}
  서론은 연구의 배경 및 필요성, 선행연구 또는 이론적 배경, 본 연구의 목적 및 내용 등으로 구성한다.

  \subsection{서론의 기능 및 모양}
  논문의 서론은 다음의 기능을 갖는다.
  \begin{enumerate}
    \item 독자의 흥미를 유발 (직접적, 특정, 짧게 그리고 명백하게)
    \item 독자가 논문을 이해할 수 있게 충분한 내용을 제공
    \item 가장 중요한 문장은 연구질문 (research questions)
  \end{enumerate}

  \begin{figure}[htbp]
      \centering
      \vspace{2cm} % 도식 배정 레이아웃 공간 확보용 가상 박스
      \caption{Shape of Introduction}
      \label{fig:intro_shape}
  \end{figure}

  \subsection{서론의 시제(tense)}
  서론의 시제는 일반적으로 다음의 규칙을 따른다.
  \begin{enumerate}
    \item 배경(일반화 된 내용)은 현재형
    \item 다른 연구자의 논문에 발표된 내용은 과거형
    \item 연구질문의 신호(signal)는 과거형 (이 연구의 목적은 \ldots 였다.)
  \end{enumerate}